An observer observed a transit of Venus near the North Pole of the Earth. The
transit path of Venus is shown in the picture below. A, B, C, D are all on the
path of transit and marking the center of the Venus disk. At A and B, the
center of Venus is superposed on the limb of the Sun disk; C corresponds to the
first contact while D to the fourth contact, $\angle AOB = 90^\circ$, MN is
parallel to AB. The first contact occured at 9:00 UT. Calculate the time of the
fourth contact.
\[ T_{\mathrm{venus}} = 224.70\ \mathrm{days}, \quad
   T_{\mathrm{earth}} = 365.25\ \mathrm{days}, \quad
   a_{\mathrm{venus}} = 0.723\,\mathrm{AU}, \quad
   r_{\mathrm{venus}} = 0.949\,r_\oplus \]
\ioaafig{2010_TH_S14-f1.pdf}
% through points A, B (centre of Venus on the solar limb), C (first contact),
% D (fourth contact), and line MN parallel to AB; printed in the 2010
% proceedings (solutions section, p. 38, left figure).
